% Front matter — Danh mục ký hiệu và chữ viết tắt (bản viết lại 2026-10-14, thay noi-dung/04-viet-tat.tex khi ghép).
% Chỉ gồm chữ viết tắt và ký hiệu có trong các chương đang ghép: Ch1 r40 v4-v4, Ch2 v6, Ch3 r33 v4-v3, Ch4A r36 v5-v3,
% Ch4B r39 v11-v2, Ch5 r31 v2, Phụ lục r31 v2. Đã bỏ các chữ tắt và ký hiệu mà GLOSSARY-VI-v2 mục H cho bỏ, cùng mọi
% mục của bản cũ không còn xuất hiện trong các chương. Chương dùng mỗi mục lần đầu (tệp:dòng): FRONT-NOTES.md mục 6–7.
\chapter*{Danh mục ký hiệu và chữ viết tắt}
\addcontentsline{toc}{chapter}{Danh mục ký hiệu và chữ viết tắt}

\vspace{4pt}
\noindent\textbf{Chữ viết tắt}
\vspace{4pt}

\begin{longtable}{@{}p{2.6cm}p{4.6cm}p{7.6cm}@{}}
\toprule
\textbf{Viết tắt} & \textbf{Cụm từ đầy đủ} & \textbf{Ý nghĩa}\\
\midrule
\endhead
CH1, CH2, CH3 & Câu hỏi nghiên cứu 1, 2, 3 & Ba câu hỏi nghiên cứu nêu ở mục~\ref{muc:c1-cauhoi}.\\
DARP & Dial-a-Ride Problem & Bài toán đưa đón theo yêu cầu.\\
JSON & JavaScript Object Notation & Định dạng văn bản của mỗi bản tin trao đổi giữa bộ mô phỏng và bộ chạy
RideBound.\\
LF & Line Feed & Kiểu ký tự xuống dòng của các tệp mã nguồn khi chạy lại kiểm thử tự động
(Phụ lục~\ref{phuluc:C}).\\
NOAA & National Oceanic and Atmospheric Administration & Cơ quan Quản lý Khí quyển và Đại dương Quốc gia Hoa Kỳ;
nguồn số liệu tuyết của trạm Central Park.\\
TLC & Taxi \& Limousine Commission & Ủy ban Taxi và Xe limousine New York, nơi công bố dữ liệu chuyến đi
taxi.\\
\bottomrule
\end{longtable}

\vspace{8pt}
\noindent\textbf{Ký hiệu}
\vspace{4pt}

\begin{longtable}{@{}p{2.6cm}p{12.2cm}@{}}
\toprule
\textbf{Ký hiệu} & \textbf{Ý nghĩa}\\
\midrule
\endhead
\multicolumn{2}{@{}l}{\textit{Lời hứa và giờ trả (tính bằng giây trên đồng hồ mô phỏng)}}\\
$r$ & Một khách đã có lời hứa; chỉ số dưới $r$ chỉ đại lượng của khách đó.\\
$v$ & Giờ trả đã công bố cho khách ở lần gần nhất.\\
$e$ & Giờ trả tính lại nếu xe giữ nguyên kế hoạch cũ, từ trạng thái và thời gian đi hiện tại.\\
$p$ & Giờ trả nếu chọn phương án đang xét, với cùng trạng thái và thời gian đi.\\
$p_0$ & Lời hứa đầu: giờ trả công bố khi khách xác nhận đặt xe, chưa cộng dự báo.\\
$p_n$; $j$ & Lời hứa giờ trả sau lần cập nhật cuối; $j = 1, \ldots, n$ đánh số các lần cập nhật lời hứa của một
khách.\\
$D$; $D_0$ & Lời hứa theo dự báo của một lần công bố (mục~\ref{muc:dubao}); $D_0$ là lời hứa đầu theo dự báo, và
$D_0 = p_0$ khi không có dự báo.\\
$a$ & Giờ trả thực tế trong mô phỏng.\\
\multicolumn{2}{@{}l}{\textit{Các phần của độ lệch giờ trả}}\\
$x$ & $x = e - v$: phần thay đổi khi giữ nguyên kế hoạch, ở một lần cập nhật.\\
$c$ & $c = p - e$: phần dời do điều phối, gọi là phần tính cho điều phối, ở một lần cập nhật.\\
$z$ & $z = p - v = x + c$: phần khách thấy ở một lần cập nhật.\\
$q$ & $q = \sum_j c_j$: tổng có dấu của phần tính cho điều phối qua cả chuyến; dời muộn dương, dời sớm âm.\\
$f$ & $f = D_0 - p_0$: phần dự báo đã cộng sẵn vào lời hứa đầu.\\
$\rho$ & $\rho = a - p_n$: độ lệch giữa giờ trả thực tế và lời hứa cuối.\\
\multicolumn{2}{@{}l}{\textit{Luật lời hứa và chỉ số}}\\
$B$ & Trần, tính bằng giây. Với luật C: tổng $|c|$ tối đa của một khách trong cả chuyến. Với luật hạn chót: số giây
giờ trả được phép muộn hơn $p_0$.\\
$L_r$ & Phần trần khách $r$ đã dùng: tổng các $|c|$ đã tính cho $r$ từ lúc xác nhận đặt xe; $B - L_r$ là phần
trần còn lại. Viết gọn $L$ khi không cần chỉ rõ khách.\\
(N) & Giả thiết của Nhận xét~\ref{nx:giunguyen}: xe dựng được phương án giữ nguyên kế hoạch, và giờ trả của phương
án đó được tính bằng đúng phép chiếu dùng cho $e$.\\
$T$ & Ngưỡng trễ trong các chỉ số, bằng 30, 60, 120, 300 hoặc 600 giây; trễ quá $T$ nghĩa là
$a - D_0 > T$ (Bảng~\ref{bang:c4-chiso}).\\
\multicolumn{2}{@{}l}{\textit{Tên cấu hình, chỉ dùng trong bảng và hình (Bảng~\ref{bang:luat})}}\\
U & Điều phối không trần: giữ xe, điểm đón, điểm trả và giờ đón đã xác nhận; không giới hạn phần dời giờ trả.\\
FREE & Ghép tự do: không khóa gì và không kiểm lời hứa.\\
C-$B$ & Luật C với trần $B$ giây, ví dụ C-30, C-120.\\
M$^{+}$-$B$ & Luật hạn chót: giờ trả không được muộn hơn lời hứa đầu $p_0$ quá $B$ giây.\\
V-$B$ & Luật trần độ dời khách thấy: tổng $|z|$ của một khách trong cả chuyến không vượt $B$.\\
\multicolumn{2}{@{}l}{\textit{Đơn vị}}\\
s & Giây; chỉ dùng trong bảng và hình.\\
\bottomrule
\end{longtable}

\noindent Các ký hiệu chỉ dùng trong một công thức, như ở công thức~\eqref{eq:c2-chon}, \eqref{eq:chon}
và~\eqref{eq:keodan}, được giải thích ngay cạnh công thức đó.

\clearpage
